% ARCHIVO GENERADO por scripts/generar_capitulos.py. No editar a mano.
\section{Requerimientos de la iteración 2}
Aprobados por el equipo el 2026-11-15 (versión 1.4) y actualizados a la versión 1.5 por la solicitud de cambio CHG-003 (2026-11-22). Las reglas RN-01 a RN-24 de la iteración 1 siguen vigentes.

\subsection{Respuestas del equipo (2026-11-15)}

\begin{longtable}{>{\raggedright\arraybackslash}p{2.4000000000000004cm}>{\raggedright\arraybackslash}p{9.9cm}>{\raggedright\arraybackslash}p{2.1cm}}
\toprule
\textbf{Pregunta} & \textbf{Respuesta} & \textbf{Dónde quedó}\\
\midrule\endhead
P-26 Permisos & El \textbf{operador} registra, modifica, inactiva y activa propietarios y vehículos. El administrador hace todo. & Matriz de permisos (sección 3)\\
P-27 Horarios & Son \textbf{fijos}: el operador \textbf{no puede iniciar sesión fuera de su turno}, salvo que el administrador lo autorice desde su sesión. & RN-31, RN-32, RF-107, RF-110\\
P-28 Celdas & Las celdas son \textbf{fijas}: no se agregan ni cambian de tipo. \textbf{Reemplazada por CHG-003 (2026-11-22):} el administrador crea y administra las celdas. & RF-103 (\emph{Must}), RN-48 a RN-53\\
P-29 Vehículos & Placa mal escrita: se \textbf{inactiva} y se crea un \textbf{registro nuevo}. Cambio de propietario: \textbf{petición al administrador}, y debe quedar el registro del propietario anterior. & RN-26, RN-33, RF-109\\
P-30 Teléfono & \textbf{Celular colombiano.} & RN-30\\
P-31 Operadores & El administrador puede editar el nombre y restablecer la contraseña (``control sobre todo el sistema''). & RF-106\\
P-32 Trazabilidad & Se debe \textbf{mostrar qué operador hizo qué operación.} & RF-111, DEC-013\\
\bottomrule
\end{longtable}




\subsection{Reglas de negocio propuestas}

\begin{longtable}{>{\raggedright\arraybackslash}p{2.6cm}>{\raggedright\arraybackslash}p{10.0cm}>{\raggedright\arraybackslash}p{1.8cm}}
\toprule
\textbf{ID} & \textbf{Regla} & \textbf{Origen}\\
\midrule\endhead
RN-25 & No se registran dos propietarios con la misma cédula ni dos vehículos con la misma placa. & [I] (ya lo garantiza el modelo)\\
RN-26 & La placa no se modifica. Si se registró mal, el vehículo se inactiva y se registra de nuevo con la placa correcta. & [U] P-29\\
RN-27 (reemplazada) & \emph{Reemplazada por CHG-003:} las celdas ya no son fijas (ver RN-48 a RN-53). & [U] P-28\\
RN-48 & Solo el \textbf{administrador} crea, modifica, inactiva, reactiva y elimina celdas. El operador solo las consulta en el mapa. & [E] CHG-003, P-52\\
RN-49 & \textbf{Código de celda:} único y con el formato de su tipo: \texttt{A-} y 3 dígitos para carro/camioneta (\texttt{A-001} a \texttt{A-999}) y \texttt{M-} y 2 dígitos para moto (\texttt{M-01} a \texttt{M-99}). Se escribe en mayúsculas; el sistema sugiere el siguiente código libre. El código no se modifica, salvo al cambiar el tipo de una celda que nunca se usó (RN-52): como el formato depende del tipo, recibe un código del tipo nuevo. & [E] CHG-003, P-53, P-55\\
RN-50 & \textbf{Celda ocupada:} no se puede inactivar, cambiar de tipo ni eliminar mientras tenga un vehículo dentro. & [E] CHG-003\\
RN-51 & \textbf{Celda inactiva} (por ejemplo, en mantenimiento): se inactiva con un motivo obligatorio; no se ofrece al registrar una entrada; en el mapa aparece marcada como inactiva. Se puede reactivar. & [U] P-54\\
RN-52 & \textbf{Modificar y eliminar:} solo se cambia el tipo de una celda libre que nunca se ha usado. Solo se elimina una celda que nunca se ha usado; si ya tuvo estadías, se inactiva (RN-07). & [U] P-52, P-55\\
RN-53 & Cada creación, modificación, inactivación, reactivación o eliminación de una celda queda en la \textbf{bitácora} (RN-34). No hay un número máximo de celdas fuera del límite del formato del código. & [U] P-56\\
RN-29 & Al inactivar un propietario o un vehículo, si está dentro puede salir, pero no puede volver a entrar mientras siga inactivo (RN-24). & [I]\\
RN-30 & El teléfono del propietario es un celular colombiano: 10 dígitos que empiezan por 3 (por ejemplo, 3001234567). & [U] P-30\\
RN-31 & Un operador solo puede iniciar sesión dentro de su turno. El administrador no tiene restricción de horario. & [U] P-27\\
RN-32 & El administrador puede autorizar a un operador para ingresar fuera de su turno, \textbf{por el tiempo que el administrador defina} (desde, hasta). La autorización queda registrada con quién la dio y el motivo. & [U] P-27, P-34\\
RN-33 & El cambio de propietario de un vehículo lo solicita un operador y lo aprueba o rechaza el administrador. Al aprobarse, el propietario anterior queda en el historial del vehículo con sus fechas. Los pagos siguen asociados al vehículo. & [U] P-29\\
RN-35 & \textbf{Fin del turno con la sesión abierta:} (a) si el operador está realizando \textbf{cualquier registro} (propietario, vehículo, entrada, salida, novedad, etc.), el sistema \textbf{no} cierra la sesión mientras siga escribiendo, y la cierra cuando termine la operación (P-38). Si pasan \textbf{5 minutos sin que escriba nada}, se aplican los 5 minutos de aviso de (b); (b) si no, el sistema le avisa y le da \textbf{5 minutos} para cerrar sesión; pasado ese tiempo, la sesión se cierra automáticamente. Una autorización vigente (RN-32) evita el cierre. & [U] P-34\\
RN-36 & Un operador \textbf{sin horario asignado no puede iniciar sesión}. Se admiten turnos \textbf{nocturnos} que cruzan la medianoche (por ejemplo, 22:00–06:00). & [U] P-35\\
RN-37 & Las franjas de un mismo operador en un mismo día \textbf{no pueden cruzarse}. & [U] (revisión de HU-020)\\
RN-34 & Toda operación que crea, modifica, activa o inactiva datos queda registrada con el usuario, la fecha y hora, y qué cambió. & [U] P-32\\
\bottomrule
\end{longtable}




\subsection{Matriz de permisos (RNF-002)}

\begin{longtable}{>{\raggedright\arraybackslash}p{10.2cm}>{\raggedright\arraybackslash}p{1.6cm}>{\raggedright\arraybackslash}p{2.6cm}}
\toprule
\textbf{Operación} & \textbf{Operador} & \textbf{Administrador}\\
\midrule\endhead
---|:-:|:-: &  & \\
Registrar, modificar, activar e inactivar propietarios y vehículos & $\checkmark$ & $\checkmark$\\
Solicitar un cambio de propietario & $\checkmark$ & $\checkmark$\\
Aprobar o rechazar un cambio de propietario & --- & $\checkmark$\\
Consultar propietarios, vehículos y el mapa de celdas & $\checkmark$ & $\checkmark$\\
Crear, modificar, inactivar, reactivar y eliminar celdas (CHG-003) & --- & $\checkmark$\\
Gestionar operadores: crear, editar, restablecer contraseña, activar e inactivar & --- & $\checkmark$\\
Definir horarios y autorizar ingresos fuera de turno & --- & $\checkmark$\\
Consultar la bitácora de un vehículo o propietario (en su ficha) & $\checkmark$ & $\checkmark$\\
Consultar la bitácora completa & --- & $\checkmark$\\
\bottomrule
\end{longtable}




\subsection{Requerimientos funcionales}

\subsubsection{RF-101 — Registrar propietario}
\begin{itemize}
  \item \textbf{Actor:} Operador, Administrador.
  \item \textbf{Entradas:} nombre completo, cédula, celular, correo y dirección. Todos obligatorios.
  \item \textbf{Validaciones:} cédula de 5 a 10 dígitos y sin repetirse (RN-25); celular según RN-30; correo válido.
  \item \textbf{Prioridad:} Must.
  \item \textbf{Criterios:}
  \begin{itemize}\item Registro correcto: el propietario queda activo y la bitácora muestra quién lo registró.\end{itemize}
  \begin{itemize}\item Cédula repetida $\rightarrow$ ``Ya existe un propietario con esa cédula''.\end{itemize}
  \begin{itemize}\item Celular inválido $\rightarrow$ ``El celular debe tener 10 dígitos y empezar por 3''.\end{itemize}
\end{itemize}

\subsubsection{RF-102 — Registrar vehículo}
\begin{itemize}
  \item \textbf{Actor:} Operador, Administrador.
  \item \textbf{Entradas:} placa, tipo, marca, nombre (línea), modelo (año), color y propietario, que se elige buscando por cédula o nombre.
  \item \textbf{Validaciones:} formato de placa acorde al tipo (RN-23); placa única (RN-25); modelo entre 1950 y el año siguiente al actual; propietario activo.
  \item \textbf{Prioridad:} Must.
  \item \textbf{Nota:} para entrar, el vehículo necesita además una mensualidad pagada (RN-03). Los pagos se registran en la iteración 3.
\end{itemize}

\subsubsection{RF-103 — Gestionar celdas (CHG-003)}
\begin{itemize}
  \item \textbf{Actor:} Administrador.
  \item \textbf{Descripción:} crear una celda (tipo y código, con el siguiente código sugerido), cambiar su tipo, inactivarla con motivo, reactivarla y, si nunca se usó, eliminarla. Ver la lista de celdas con su estado (libre, ocupada o inactiva), filtrar por tipo y estado, y buscar por código.
  \item \textbf{Validaciones:} código único y con el formato de su tipo (RN-49); una celda ocupada no se inactiva, no cambia de tipo ni se elimina (RN-50); solo se cambia el tipo o se elimina una celda que nunca se usó (RN-52).
  \item \textbf{Reglas:} RN-48 a RN-53.
  \item \textbf{Prioridad:} Must.
  \item \textbf{Historial:} en la versión 1.1 quedó como \emph{Won't} porque las celdas eran fijas (P-28). La solicitud de cambio CHG-003 lo incorporó después de la observación del profesor.
\end{itemize}

\subsubsection{RF-104 — Mapa de celdas}
\begin{itemize}
  \item \textbf{Actor:} Operador, Administrador.
  \item \textbf{Descripción:} todas las celdas por zona (A: carro/camioneta, M: moto), con su estado (libre, ocupada con la placa, o inactiva) y contadores. El inventario es el que defina el administrador (RF-103).
  \item \textbf{Prioridad:} Must.
\end{itemize}

\subsubsection{RF-105 — Modificar, activar e inactivar propietarios y vehículos}
\begin{itemize}
  \item \textbf{Actor:} Operador, Administrador.
  \item \textbf{Descripción:} corregir datos y cambiar el estado activo/inactivo. Nunca se borra (RN-07).
  \item \textbf{Restricciones:} la placa no se edita (RN-26); el propietario no se cambia aquí (RF-109).
  \item \textbf{Reactivación tras el pago:} pasa a la iteración 3 (P-33). En la iteración 2, la activación e inactivación son manuales.
  \item \textbf{Prioridad:} Must.
\end{itemize}

\subsubsection{RF-106 — Gestionar operadores (completar)}
\begin{itemize}
  \item \textbf{Actor:} Administrador.
  \item \textbf{Agrega a HU-010:} editar el nombre y restablecer la contraseña.
  \item \textbf{Prioridad:} Must.
\end{itemize}

\subsubsection{RF-107 — Horarios de operador}
\begin{itemize}
  \item \textbf{Actor:} Administrador.
  \item \textbf{Descripción:} define el turno de cada operador. Al iniciar sesión, el sistema verifica que la hora actual esté dentro del turno (RN-31); si no lo está: ``Estás fuera de tu turno. Solicita autorización al administrador''. El encabezado muestra las horas del turno, por ejemplo ``Operador · Turno 06:00 – 14:00'' (sin nombre de turno).
  \item \textbf{Turnos:} franjas por día de la semana; pueden cruzar la medianoche (RN-36). Sin horario no hay ingreso (RN-36). Al terminar el turno se aplica RN-35. Puede haber \textbf{varias franjas el mismo día}, si el administrador las asigna (P-37).
  \item \textbf{Prioridad:} Must.
\end{itemize}

\subsubsection{RF-108 — Consultar propietarios y vehículos}
\begin{itemize}
  \item \textbf{Actor:} Operador, Administrador.
  \item \textbf{Descripción:} buscar por nombre, cédula o placa. La ficha del propietario muestra sus vehículos, si están dentro o fuera y su estado. La ficha del vehículo muestra el \textbf{historial de propietarios} (RN-33) y la bitácora de cambios (RF-111).
  \item \textbf{Prioridad:} Must.
\end{itemize}

\subsubsection{RF-109 — Solicitar y aprobar un cambio de propietario (nuevo)}
\begin{itemize}
  \item \textbf{Actor:} Operador (solicita) y Administrador (aprueba o rechaza).
  \item \textbf{Flujo:}
  \begin{itemize}\item El operador elige el vehículo y el nuevo propietario, que debe estar registrado, y escribe el motivo.\end{itemize}
  \begin{itemize}\item La solicitud queda \textbf{pendiente}.\end{itemize}
  \begin{itemize}\item El administrador la ve en una bandeja y la aprueba o rechaza con una observación.\end{itemize}
  \begin{itemize}\item Si la aprueba, el vehículo pasa al nuevo propietario y el anterior queda en el historial con las fechas.\end{itemize}
  \item \textbf{Reglas:} RN-33. Solo puede haber una solicitud pendiente por vehículo [I].
  \item \textbf{Prioridad:} Must.
\end{itemize}

\subsubsection{RF-110 — Autorizar ingreso fuera de turno (nuevo)}
\begin{itemize}
  \item \textbf{Actor:} Administrador.
  \item \textbf{Descripción:} desde su sesión, el administrador autoriza a un operador a ingresar fuera de su turno, con un motivo y una vigencia que él define (desde y hasta). Mientras esté vigente, el operador puede iniciar sesión y no se le cierra la sesión por fin de turno (RN-32, RN-35). Queda en la bitácora.
  \item \textbf{Prioridad:} Must.
\end{itemize}

\subsubsection{RF-111 — Bitácora de operaciones (nuevo)}
\begin{itemize}
  \item \textbf{Actor:} Administrador (bitácora completa) y Operador (solo la de un vehículo o propietario, en su ficha; P-36).
  \item \textbf{Descripción:} registro automático de cada operación (registrar, modificar, activar, inactivar, solicitar o aprobar un cambio de propietario, crear o editar operadores, horarios, autorizaciones). Muestra quién, cuándo, qué entidad y qué cambió. Se puede filtrar por usuario, fecha y tipo de operación.
  \item \textbf{Reglas:} RN-34.
  \item \textbf{Prioridad:} Must.
  \item \textbf{Implementación:} bitácora de operaciones (DEC-013, aprobada).
\end{itemize}


